\chapter{Introduction}
\label{chap:introduction}

% ============================================================
\section{Background and Motivation}
% ============================================================

Information and Communication Technology (ICT) has become a foundational
pillar of the global digital economy, supporting mobile connectivity, cloud
services, digital governance, and other critical infrastructure. As of 2025,
mobile networks serve approximately 5.7 billion unique subscribers
worldwide, representing over 70\% of the global population
\citep{gsma2025_mobile_economy}. Continued growth in mobile data demand and
the expansion of 4G and 5G networks are placing increasing pressure on the
energy consumption of telecom infrastructure.

At the core of this infrastructure lies the \textbf{telecom tower site} (or
base transceiver station, BTS site), which must operate continuously to
maintain uninterrupted service. A typical multi-technology tower consumes
2--8\,kW continuously, making the Radio Access Network (RAN) responsible for
60--80\% of a mobile operator's total electricity demand
\citep{americas2023_energy_efficiency}. In developing regions such as India
and Sub-Saharan Africa, unreliable grid access means that many towers rely
heavily on diesel generators and battery banks to maintain service continuity
\citep{gsma2020_tower_esco}. Diesel fuel represents 25--30\% of tower OPEX,
rising to 70\% at remote sites \citep{datoms2024,gsma2020_tower_esco}.

Fuel logistics failures---through pilferage, stockouts, and delivery
disruptions---pose significant risks to tower reliability; the 2025 threat
to 16{,}000 Nigerian sites underscores diesel supply-chain vulnerability
\citep{datoms2024,alton2025}. In weak-grid regions, the resulting reliance
on diesel generators drives both high operating costs and substantial carbon
emissions. In the Indian context, a typical macro BTS site operating a
diesel generator for approximately eight hours per day consumes about
8{,}760\,L of diesel annually, corresponding to approximately
23--24\,t CO\textsubscript{2} per site per year when using a standard
emission factor of 2.68\,kg\,CO\textsubscript{2} per litre of diesel
\citep{trai2011_green_telecom,ghgprotocol2017_emission_factors}.

Global sustainability commitments further increase the need to reduce
diesel dependence. The GSMA Climate Action roadmaps and major operators
target net-zero emissions by 2040--2050, while diesel-powered sites remain
one of the largest controllable sources of Scope\,1 emissions
\citep{gsma2023}. Hybrid renewable systems (solar PV + battery + diesel)
have demonstrated 20--65\% reductions in diesel consumption and
CO\textsubscript{2} emissions in field deployments
\citep{gsma2020_tower_esco,datoms2024}. However, maximising these savings
under stochastic solar irradiance, traffic-driven loads, and grid outages
requires more sophisticated control than current rule-based systems.

In India---the world's second-largest mobile market with over
1.2 billion subscribers and over
800{,}000 tower sites \citep{dot2024_towers}---the challenge is particularly
acute. Rural and semi-urban grid availability often falls below 14 hours per
day, making diesel dominant at many sites \citep{trai2024,datoms2024}.
Aggressive 5G rollout and renewable energy mandates intensify pressure for
low-carbon hybrid transitions while controlling OPEX.


% ============================================================
\section{Current Practices and Their Limitations}
% ============================================================

Industry-standard energy management today relies predominantly on
\textbf{rule-based} or \textbf{threshold-driven controllers} embedded in
Power Management Units (PMUs) or Hybrid Power Controllers (HPCs). Typical
rules include: use solar whenever available; switch to battery when the grid
fails; and start the diesel generator if battery State-of-Charge (SoC) falls
below 30\%.

These heuristics are simple, robust, and inexpensive, but they are
fundamentally \emph{myopic}: they ignore long-term fuel cost, battery
degradation cost, forecasted outages, variable diesel prices, and potential
carbon penalties. Tools such as \textbf{HOMER Pro} are widely used for system
design and planning. They can produce static daily dispatch schedules and size
PV/DG/battery systems that reduce diesel use by 30--58\%, yet they operate
offline, typically under deterministic or average conditions, and cannot adapt
in real time to uncertainty \citep{homer_energy_2020,meena2020_homer_telecom}.

More advanced approaches---such as Model Predictive Control (MPC), stochastic
programming, and classical optimisation---improve performance by explicitly
optimising over a rolling horizon. MPC and related methods require an
explicit system model, suitable forecasts, and repeated online
optimisation; their effectiveness therefore depends on model and forecast
quality as well as the available controller resources.

Reinforcement Learning (RL) has recently shown strong results in microgrids,
electric vehicles, and building energy management, where agents learn to trade
off cost and reliability under uncertainty
\citep{mao2021_inventory_rl,lin_ev_rl_2024,wang2022_building_rl}. Recent
work specifically on telecom base station energy management has demonstrated
18--47\% reductions in operating cost relative to rule-based and MPC baselines
\citep{ma_pan_2025_telecom_ies,yan_2025_safe_marl_bs,alhajhassan_2024_battery_ageing_rl}.
However, an important modelling gap remains across this body of work: diesel fuel is
generally treated as an always-available, unconstrained resource. To the best
of the author's knowledge, no prior RL formulation for telecom towers models
the diesel tank level as a finite depletable stock or treats replenishment
ordering as a control action. Existing formulations also do not account for
stochastic delivery lead times and their interaction with dispatch decisions.
Most existing approaches optimise dispatch and fuel replenishment
independently, despite the operational coupling between diesel consumption
and future inventory availability under uncertain delivery lead times. This
gap motivates the
methodology adopted in this thesis. The following chapter reviews related
work across telecom energy management, microgrid scheduling, and
inventory-aware control, and positions the proposed formulation relative
to existing approaches.


% ============================================================
\section{Research Problem and Objectives}
% ============================================================

The problem addressed here is how to coordinate energy resources
and diesel replenishment in hybrid telecom towers under significant uncertainty
in solar generation, grid availability, and diesel delivery, so as to minimise
unserved energy and operational expenditure while improving service
reliability and respecting operational constraints.

The central research question is:

\begin{quote}
\emph{Does jointly learning diesel ordering and energy dispatch within a single
reinforcement-learning agent reduce unserved energy at remote telecom
base stations relative to approaches that treat these decisions separately ---
including both classical rule-based decomposition and optimisation-based
dispatch?}
\end{quote}

To answer this question, the thesis formulates the problem as a
\textbf{Constrained Markov Decision Process (CMDP)} and evaluates three
hypotheses. The empirical evaluation focuses on reliability
measured through \textbf{Expected Energy Not Served (EENS)} and related
operational metrics, and deliberately includes a Model Predictive Control
(MPC) baseline and a perfect-information Oracle-MPC reference alongside the
reinforcement-learning comparators, so that the effects of replenishment
learning, dispatch quality, forecasting, and action masking can be
examined separately:

\begin{description}
  \item[H1 (Joint ordering).] Does joint learning of dispatch and diesel
  ordering reduce unserved energy relative to a classical $(s,S)$ ordering
  policy, when dispatch policies remain broadly comparable despite the
  reduced action space?
\item[H2 (Action masking).] Does hard action masking improve dispatch-policy
robustness under stressed logistics, and does its effect depend on whether
diesel replenishment is externally managed or learned jointly?
  \item[H3 (Logistics-observation hypothesis).] Does removing the five inventory- and
  delivery-related observation channels reduce reliability relative to the
  full RLInv observation space?
\end{description}

The specific objectives of this work are to:

\begin{enumerate}
  \item Formulate the hybrid telecom tower energy management problem as a CMDP
  that treats battery state-of-charge and diesel fuel level as coupled
  state variables under uncertain supply and demand, with diesel
  modelled as a \textbf{finite replenishable inventory} subject to stochastic
  delivery lead times.

  \item Design and train reinforcement learning agents capable of jointly
  learning energy dispatch and diesel replenishment strategies under
  operational constraints, alongside decomposed control baselines in
  which ordering and dispatch are learned or optimised separately.

  \item Develop a Gymnasium-compatible simulation environment driven by real
  ITU/Zindi telecom site traces that reproduces empirically grounded load,
  PV generation, and grid availability dynamics at hourly resolution.

  \item Benchmark the learned RL policies against rule-based,
  decomposition-based, and optimisation-based baselines --- including MPC and
  a perfect-information Oracle-MPC reference --- across ten sites and four
  delivery-logistics scenarios of increasing severity, with the learned
  policies trained across ten independent seeds. Quantify improvements in
  unserved energy, diesel
  consumption, and operational reliability, with paired seed-level
  statistical analysis for the core policy and ablation comparisons.
\end{enumerate}


% ============================================================
\section{Contributions}
% ============================================================

The main contributions of this thesis are:

\begin{enumerate}
  \item A telecom-tower-specific, inventory-theoretic CMDP formulation that
  explicitly represents battery state-of-charge and diesel fuel volume as
  state variables subject to uncertain demand, stochastic replenishment lead
  times, and operational cost, reliability, and feasibility penalties.

  \item The design and empirical evaluation of a constrained reinforcement
  learning framework---\textbf{RLInv}---for hybrid telecom tower energy
  management, including inventory-aware state and action design with hard
action masking to prevent infeasible control actions. RLInv is benchmarked against
  a dispatch-only RL policy with externally
  managed replenishment (TrackB), optimisation-based dispatch policies
  (MPC and Oracle-MPC), and the A6 ablation. The RLInv--A6
  comparison provides a cleaner practical replenishment comparison than
  TrackB, while the reduced dispatch action space used by A6 remains an
  acknowledged design difference, not present in
  the RL-based telecom energy management literature reviewed in
  Chapter~\ref{chap:literature}.

  \item A reproducible simulation environment and experimental protocol
  integrating the ITU~2024 Smart Energy Scheduling dataset (ten base station
  sites, 60~days each) with a synthetic lead-time overlay, evaluated across
  four delivery-logistics scenarios with the learned policies trained across
  ten independent seeds, with
  paired seed-level statistical analysis for the core policy and ablation
  comparisons.

  \item A scenario-dependent empirical characterisation of when joint
  optimisation is beneficial: the H1 comparison favours RLInv under the Normal
and Extreme logistics scenarios, while no statistically distinguishable advantage
  is observed under Delayed or Monsoon conditions.
  
   \item A larger architecture-dependent gap is observed under stressed
  logistics in the externally managed replenishment comparison
  (A6 vs.\ TrackB), while the cleaner RLInv--A7 comparison finds no
  measurable independent mean-EENS benefit from masking when
  replenishment is already learned jointly.

  \item An exploratory extension investigating explicit delivery-delay
  tracking (RLInv-DA) is also reported: explicit
  delay awareness did not provide a consistent improvement over RLInv under
  the tested conditions.
\end{enumerate}


% ============================================================
\section{Scope and Limitations}
% ============================================================

To maintain focus and analytical tractability, this thesis makes the following
scope choices:

\begin{itemize}
  \item The control problem is formulated at the level of an individual tower
  site; empirical evaluation is conducted across all ten ITU/Zindi sites to
  assess performance across heterogeneous site profiles.
  \item Control decisions are taken at \textbf{hourly resolution}
  ($\Delta t = 1$\,h), consistent with the temporal granularity of the
  ITU/Zindi dataset.
  \item Component sizes (PV, battery, DG rating, tank capacity) are assumed
  given; the focus is on \emph{operational control}, not investment planning.
  \item Battery operation is modelled using engineering constraints on
  state of charge and energy balance; detailed electrochemical
  degradation is outside the scope of this work.
  \item The stochastic diesel-delivery process is the only synthetic stochastic
process introduced into the principal experiments. Load, PV generation, and grid-availability inputs
  are drawn directly from the ITU site traces, while component parameters
  and operational constraints are specified from the adopted engineering
  assumptions.
  \item Statistical comparisons use $n = 10$ independent seeds with paired
  testing across identical site--scenario combinations (an increase from
  an $n=3$ pilot evaluation conducted earlier in the project). This
  improves statistical power considerably, but for some comparisons ---
  particularly the inventory-observability ablation (H3) --- the
  differences remain small and are not statistically distinguishable
  under the final protocol; additional replication could clarify whether
  a small effect exists.
\end{itemize}

Limitations arising from these choices include:

\begin{itemize}
  \item The proposed framework is trained under representative
  delivery-delay regimes defined in the experimental protocol and
  subsequently evaluated across four representative logistics scenarios.
  Chapter~\ref{chap:methodology} describes the training protocol in
  detail, while Chapter~\ref{chap:results} presents the corresponding
  evaluation results. RLInv's advantage over classical ordering is
  scenario-dependent rather than uniform, and narrows at very long delivery
  lead times as physical tank capacity becomes the binding constraint for
  every policy evaluated, not a training-distribution failure specific to
  RLInv (Chapter~\ref{chap:results}).
  \item Multi-site coordination and shared fuel logistics are not modelled.
  \item Integrated forecasting modules for load or solar are not learned jointly
  with the control policy.
  \item The diesel generator is modelled at fixed rated power when ON; continuous
  power setpoint control is left for future work.
\end{itemize}

These aspects are identified as promising directions for future work.
% ============================================================
\section{Thesis Roadmap and Relation Between Chapters}
% ============================================================

To make the logical flow of the thesis explicit,
Fig.~\ref{fig:thesis-roadmap} summarises how the chapters relate to
one another and highlights where the core technical contributions arise. The
thesis progresses from motivation and gap identification, through formal
problem definition and implementation, to empirical evaluation and
conclusions. The main contribution lies in connecting three elements
that are typically treated separately in prior literature:
(i)~telecom tower hybrid-energy dispatch,
(ii)~diesel inventory and replenishment uncertainty, and
(iii)~constrained reinforcement learning for joint operational control.
Chapter~3 formulates the problem, Chapter~4 presents the proposed method,
and Chapter~7 reports the empirical evidence; Chapters~3--7 together form
the core technical workflow of the thesis.

\begin{figure}[H]
\centering
\resizebox{0.82\textwidth}{!}{%
\begin{tikzpicture}[
    node distance=0.4cm and 1.6cm,
    every node/.style={font=\small, align=center},
    box/.style={
        draw, rounded corners, rectangle,
        minimum width=3.4cm, minimum height=1.05cm
    },
    novel/.style={
        draw, rounded corners, rectangle,
        minimum width=3.4cm, minimum height=1.05cm,
        fill=gray!15, thick
    },
    arrow/.style={->, thick},
    xarrow/.style={->, thick, dashed}
]
\node[box]   (c1) {Chapter 1\\Motivation, scope,\\research objectives};
\node[box,   below=of c1] (c2) {Chapter 2\\Literature review\\and research gap};
\node[novel, below=of c2] (c3) {\textbf{Chapter 3}\\CMDP: joint dispatch\\+ inventory formulation};
\node[novel, below=of c3] (c4) {\textbf{Chapter 4}\\Methodology: TelecomEnv,\\RLInv, MaskablePPO};
\node[box,   below=of c4] (c5) {Chapter 5\\Dataset, EDA, site\\difficulty classification};
\node[box,   below=of c5] (c6) {Chapter 6\\Experimental execution\\and training runs};
\node[novel, below=of c6] (c7) {\textbf{Chapter 7}\\Results, ablations,\\robustness analysis};
\node[box,   below=of c7] (c8) {Chapter 8\\Conclusion, limitations,\\future work};

\draw[arrow] (c1) -- (c2);
\draw[arrow] (c2) -- (c3);
\draw[arrow] (c3) -- (c4);
\draw[arrow] (c4) -- (c5);
\draw[arrow] (c5) -- (c6);
\draw[arrow] (c6) -- (c7);
\draw[arrow] (c7) -- (c8);

\draw[xarrow] (c3.east) to[out=0, in=0, looseness=1.4]
    node[right, font=\footnotesize] {CMDP defines\\eval metrics}
    (c7.east);
\draw[xarrow] (c5.east) to[out=0, in=0, looseness=0.9]
    node[right, font=\footnotesize] {ITU data\\grounds results}
    (c7.east);
\end{tikzpicture}%
}
\caption{Thesis roadmap. Shaded boxes (Chapters~3, 4, and~7) mark the
formulation, proposed method, and empirical evidence that together form
the core technical contribution; Chapters~3--7 as a whole constitute the
core technical workflow. Dashed arrows indicate cross-chapter
dependencies. The vertical spine shows the sequential reading order.}
\label{fig:thesis-roadmap}
\end{figure}
\FloatBarrier

% ============================================================
\section{Thesis Organisation}
% ============================================================

Following the roadmap shown in Fig.~\ref{fig:thesis-roadmap}, the
remainder of this thesis elaborates each stage of the research workflow
in detail:

\begin{itemize}
  \item \textbf{Chapter~\ref{chap:literature}} reviews telecom, microgrid,
  and inventory-theoretic literature to identify the gap that motivates
  this work.

  \item \textbf{Chapter~\ref{chap:problem}} formulates the hybrid telecom
  tower system as a constrained Markov decision process, including the
  state and action spaces, cost function, and operational constraints.

  \item \textbf{Chapter~\ref{chap:methodology}} presents the implemented
  solution: the \texttt{TelecomEnv} simulator, the \texttt{MaskablePPO}
  agent, the full comparison-policy set, and the evaluation protocol.

  \item \textbf{Chapter~\ref{chap:datasets-setup}} documents the ITU/Zindi
  dataset and the exploratory analysis that informed simulator design and
  site selection.

  \item \textbf{Chapter~\ref{chap:experimental-execution}} records how the
  experimental plan was executed, from training runs through convergence
  monitoring to the evaluation campaign.

  \item \textbf{Chapter~\ref{chap:results}} presents and interprets the
  full experimental results, including the primary comparison, the H1--H3
  ablations, and the robustness checks.

  \item \textbf{Chapter~\ref{chap:conclusion}} discusses what the results
  mean, summarises the contributions and limitations, and outlines
  directions for future work.
\end{itemize}

\clearpage